\section{Polar crystals}\label{app:D.14}

In a polar crystal the field of a longitudinal optical wave is long-ranged (\cref{sec:5.5}). It couples the longitudinal phonon to an electron with a strength that diverges as $1/Q$, the Fr\"ohlich vertex~\cite{frohlich1954electrons}. This section derives the Lyddane--Sachs--Teller relation~\cite{lyddane1941polar} for several polar modes and the Fr\"ohlich vertex from the energy function \cref{eq:5.5-U} of the displacements and the field, in the form \cref{sec:5.5,sec:6.2} use. The approach is that of Huang's model as presented by Born and Huang~\cite{born1954dynamical} and Baroni et al.~\cite[Sect.~II.D.2]{baroni2001phonons}, generalized to anisotropic tensors~\cite{cochran1962dielectric,gonze1997dynamical}.

\subsection*{The Lyddane--Sachs--Teller relation}

With several polar modes the steps that led to \cref{eq:5.5-eps} give the general form. Let $\omega_{\mathrm{T}\nu}$ and $\vect{e}_{\kappa\nu}$ be the frequencies and polarization vectors of the analytic problem at $\vect{Q} = 0$, the eigenpairs of $C^{\mathrm{an}}_{\kappa\kappa'}/(M_\kappa M_{\kappa'})^{1/2}$, chosen real. Expanding $\vect{u}_\kappa = \sum_\nu\vect{e}_{\kappa\nu}a_\nu/\sqrt{M_\kappa}$, each normal coordinate obeys $\ddot a_\nu = -\omega_{\mathrm{T}\nu}^2a_\nu + e\,\vect{z}_\nu\cdot\vect{\mathcal{E}}$, with the mode charge vector $\vect{z}_\nu = \sum_\kappa\vect{e}_{\kappa\nu}\cdot Z^*_\kappa/\sqrt{M_\kappa}$, and it contributes $(e/\Omega)\,a_\nu\vect{z}_\nu$ to the polarization. In a cubic crystal, along any direction $\hat{\vect{e}}$,
\begin{equation*}
  \varepsilon(\omega) = \varepsilon_\infty + \frac{e^2}{\Omega}\sum_\nu\frac{(\vect{z}_\nu\cdot\hat{\vect{e}})^2}{\omega_{\mathrm{T}\nu}^2 - \omega^2} .
\end{equation*}
This is a rational function of $\omega^2$ that tends to $\varepsilon_\infty$ at large $\omega$, has simple poles at the polar transverse frequencies, and, by the argument of \cref{sec:5.5}, vanishes at the longitudinal frequencies $\omega_{\mathrm{L}\nu}$. Between two consecutive poles it rises monotonically from $-\infty$ to $+\infty$, and above the highest pole it rises from $-\infty$ toward $\varepsilon_\infty$, so there is exactly one zero above each pole and below the next. A rational function is fixed by its poles, its zeros and its limit at infinity, so $\varepsilon(\omega) = \varepsilon_\infty\prod_\nu(\omega_{\mathrm{L}\nu}^2 - \omega^2)/(\omega_{\mathrm{T}\nu}^2 - \omega^2)$, and at $\omega = 0$
\begin{equation}
  \frac{\varepsilon_0}{\varepsilon_\infty} = \prod_\nu\frac{\omega_{\mathrm{L}\nu}^2}{\omega_{\mathrm{T}\nu}^2} ,
  \label{eq:D.14-LSTgen}
\end{equation}
the generalized relation of Cochran and Cowley~\cite{cochran1962dielectric}.

\subsection*{The Fr\"ohlich vertex}

The field of a longitudinal phonon acts on the electrons as well. Take one mode $(\vect{Q},\nu)$ at small $\vect{Q}$ in the expansion \cref{eq:A.6-u}, $\vect{u}_{\vect{R}\kappa} = (2NM_\kappa\omega_{\vect{Q}\nu})^{-1/2}\,\vect{e}_{\kappa\nu}(\vect{Q})\,\ee^{\ii\vect{Q}\cdot\vect{R}}B_{\vect{Q}\nu}$, and replace the cell position $\vect{R}$ by a continuous $\vect{r}$, which is allowed for $|\vect{Q}|$ small compared with the reciprocal-lattice vectors. The condition $\hat{\vect{Q}}\cdot\vect{D} = 0$ of \cref{sec:5.5} gives the field of the wave,
\begin{equation*}
  \vect{\mathcal{E}}(\vect{r}) = -\hat{\vect{Q}}\;\frac{e}{\Omega}\sum_\kappa\frac{\vect{e}_{\kappa\nu}(\vect{Q})\cdot Z^*_\kappa\cdot\hat{\vect{Q}}}{(2NM_\kappa\omega_{\vect{Q}\nu})^{1/2}\;\hat{\vect{Q}}\cdot\varepsilon_\infty\cdot\hat{\vect{Q}}}\;\ee^{\ii\vect{Q}\cdot\vect{r}}B_{\vect{Q}\nu} .
\end{equation*}
An electron, with charge $-e$, feels the force $-e\vect{\mathcal{E}}$, and its potential energy $V^{\mathrm{L}}(\vect{r})$ obeys $\nabla V^{\mathrm{L}} = e\vect{\mathcal{E}}$. For a wave $\propto\ee^{\ii\vect{Q}\cdot\vect{r}}$ the gradient is $\ii\vect{Q}$. The potential energy is then $V^{\mathrm{L}} = e\,\hat{\vect{Q}}\cdot\vect{\mathcal{E}}/(\ii Q)$, or
\begin{equation*}
  V^{\mathrm{L}}_{\vect{Q}\nu}(\vect{r}) = \ii\,\frac{e^2}{\Omega}\sum_\kappa\frac{\vect{e}_{\kappa\nu}(\vect{Q})\cdot Z^*_\kappa\cdot\vect{Q}}{(2NM_\kappa\omega_{\vect{Q}\nu})^{1/2}\;\vect{Q}\cdot\varepsilon_\infty\cdot\vect{Q}}\;\ee^{\ii\vect{Q}\cdot\vect{r}}B_{\vect{Q}\nu} .
\end{equation*}
Summed over modes and written in second quantization over Bloch states with \cref{eq:2.3-onebody}, it has the form \cref{eq:A.6-Hep}, with the matrix element of $\ee^{\ii\vect{Q}\cdot\vect{r}}$ the density vertex $\rho_{mn}(\vect{k},\vect{Q})$ of \cref{eq:A.5-rho}. Reading off the factor of $N^{-1/2}c^\dagger_{m,\vect{k}+\vect{Q}}c_{n\vect{k}}B_{\vect{Q}\nu}$, the long-range vertex is
\begin{equation}
  g^{\mathrm{L}}_{mn\nu}(\vect{k},\vect{Q}) = \ii\,\frac{e^2}{\Omega}\sum_\kappa\frac{\vect{e}_{\kappa\nu}(\vect{Q})\cdot Z^*_\kappa\cdot\vect{Q}}{(2M_\kappa\omega_{\vect{Q}\nu})^{1/2}\;\vect{Q}\cdot\varepsilon_\infty\cdot\vect{Q}}\;\rho_{mn}(\vect{k},\vect{Q}) .
  \label{eq:D.14-gL}
\end{equation}
It is the $\vect{G} = 0$ term of Ref.~\cite[Eq.~(4)]{verdi2015frohlich}, with $e^2_{\mathrm{SI}}/\varepsilon_0^{\mathrm{SI}} \to e^2$, the factor $N^{-1/2}$ moved into \cref{eq:A.6-Hep} and $Z^*$ transposed to our index order. The terms with $\vect{G} \ne 0$ of that equation come from the periodic arrangement of the dipoles and are regular at $\vect{Q} \to 0$. \Cref{eq:D.14-gL} obeys the Hermiticity condition \cref{eq:A.6-herm}. Under $\vect{Q} \to -\vect{Q}$ and $(m,\vect{k}) \leftrightarrow (n,\vect{k}+\vect{Q})$, the polarization vector becomes $\vect{e}_{\kappa\nu}(\vect{Q})^*$ by \cref{eq:A.6-e}, the factor $\vect{Q}$ in the numerator changes sign, and $\rho_{nm}(\vect{k}+\vect{Q},-\vect{Q}) = \rho_{mn}(\vect{k},\vect{Q})^*$; the product is the complex conjugate of \cref{eq:D.14-gL}, since $\ii \to -\ii$ is supplied by the sign of $\vect{Q}$. Only longitudinal modes couple, through $Z^*_\kappa\cdot\vect{Q}$, and the vertex grows as $1/Q$ at small $Q$.

For the cubic two-atom crystal the longitudinal optical mode at small $\vect{Q}$ has $\vect{e}_{1} = (M_2/M_{\mathrm{cell}})^{1/2}\hat{\vect{Q}}$ and $\vect{e}_{2} = -(M_1/M_{\mathrm{cell}})^{1/2}\hat{\vect{Q}}$, which satisfy \cref{eq:A.6-e}. The sum in \cref{eq:D.14-gL} is then $Z^*Q\,[(M_2/M_{\mathrm{cell}}M_1)^{1/2} + (M_1/M_{\mathrm{cell}}M_2)^{1/2}] = Z^*Q/\sqrt{\mu_{\mathrm{red}}}$, and
\begin{equation*}
  g^{\mathrm{L}} = \ii\,\frac{e^2Z^*}{\Omega\,\varepsilon_\infty\,Q\,(2\omega_{\mathrm{L}}\mu_{\mathrm{red}})^{1/2}}\;\rho_{mn}(\vect{k},\vect{Q}) .
\end{equation*}
By \cref{eq:5.5-LOTO,eq:5.5-LST}, $e^2Z^{*2}/(\Omega\varepsilon_\infty\mu_{\mathrm{red}}) = \omega_{\mathrm{L}}^2 - \omega_{\mathrm{T}}^2 = \omega_{\mathrm{L}}^2(1 - \varepsilon_\infty/\varepsilon_0)$, and the Born charge and the masses drop out:
\begin{equation}
  g^{\mathrm{L}}_{mn}(\vect{k},\vect{Q}) = \frac{\ii}{Q}\left[\frac{e^2\omega_{\mathrm{L}}}{2\Omega}\Bigl(\frac{1}{\varepsilon_\infty} - \frac{1}{\varepsilon_0}\Bigr)\right]^{1/2}\rho_{mn}(\vect{k},\vect{Q}) ,
  \label{eq:D.14-frohlich}
\end{equation}
for $Z^* > 0$, with an overall minus sign otherwise. This is Fr\"ohlich's vertex~\cites{frohlich1954electrons}[Eq.~(1)]{verdi2015frohlich}, with the unit and index changes listed above, and with the density vertex in place of the plane-wave overlap $\delta_{mn}$ of a single parabolic band. The combination $1/\varepsilon_\infty - 1/\varepsilon_0$ is the part of the screening due to the ions. The electrons see the field of the displaced ions screened by $\varepsilon_\infty$; a static charge would be screened by $\varepsilon_0$, and the difference is the field that the lattice polarization adds. In a nonpolar crystal $\varepsilon_0 = \varepsilon_\infty$, $Z^* = 0$, and the vertex vanishes. \Cref{sec:6.2} writes \cref{eq:D.14-frohlich} through Fr\"ohlich's coupling constant and evaluates it for GaAs.
